\documentclass[aspectratio=169]{beamer}
\usetheme{Madrid}
\usecolortheme{seagull}
\setbeamertemplate{navigation symbols}{}
\setbeamertemplate{footline}[frame number]
\usepackage[T1]{fontenc}
\usepackage[utf8]{inputenc}
\usepackage{lmodern}
\usepackage{graphicx}
\usepackage{amsmath,amssymb}
\usepackage{booktabs}
\usepackage{xcolor}
\usepackage{ragged2e}
\usepackage{array}
\definecolor{TitleBlue}{RGB}{18,46,84}
\definecolor{SoftGray}{RGB}{245,246,248}
\setbeamercolor{title}{fg=TitleBlue}
\setbeamercolor{frametitle}{fg=white,bg=TitleBlue}
\setbeamercolor{block title}{fg=white,bg=TitleBlue}
\setbeamercolor{block body}{bg=SoftGray}
\setbeamercolor{alertblock title}{fg=white,bg=black}
\setbeamercolor{alertblock body}{bg=SoftGray}
\setbeamertemplate{items}[square]
\setbeamersize{text margin left=0.8cm,text margin right=0.8cm}
\graphicspath{{exhibits/}}
\newcommand{\E}{\mathbb{E}}
\newcommand{\Var}{\operatorname{Var}}
\newcommand{\Cov}{\operatorname{Cov}}
\newcommand{\Ncdf}{N}
\begin{document}
\begin{frame}[t]{Global Financial Management Strategy}
\centering
\vspace{0.62cm}
{\Large\textbf{Week 9: The Black--Scholes Model from First Principles}}\\[0.5cm]
{\large Connecting Global Strategy with Financial Decisions}\\[0.3cm]
{\normalsize LIU YONGSHANG}\\[0.7cm]
\begin{tabular}{ll}
\textbf{Course:} & Global Financial Management Strategy\\
\textbf{Department:} & Business Administration\\
\end{tabular}
\end{frame}
\begin{frame}[t]{Learning Objectives}
\justifying
\begin{itemize}
\item Explain why continuous compounding produces the number $e$.
\item Distinguish a normal distribution from a lognormal distribution.
\item Interpret Brownian motion, $dt$, $dW$, and a stochastic differential equation.
\item Derive It\^o\textquotesingle s lemma from a second-order Taylor expansion.
\item Construct the Black--Scholes hedge and derive the pricing PDE.
\item Derive the European call formula and interpret $d_1$ and $d_2$.
\end{itemize}
\end{frame}

\begin{frame}[t]{Lecture Roadmap}
\justifying
\begin{enumerate}
\item[9.1] Exponential Growth and the Number $e$
\item[9.2] Probability Foundations
\item[9.3] Normal and Lognormal Distributions
\item[9.4] Random Walks and Brownian Motion
\item[9.5] Stochastic Differential Notation
\item[9.6] Geometric Brownian Motion
\end{enumerate}
\end{frame}

\begin{frame}[t]{Lecture Roadmap (continued)}
\justifying
\begin{enumerate}
\item[9.7] It\^o\textquotesingle s Lemma
\item[9.8] No-Arbitrage and Replication
\item[9.9] The Black--Scholes PDE
\item[9.10] Solving the PDE
\item[9.11] The Black--Scholes Formula
\item[9.12] Interpretation, Extensions, and Limits
\end{enumerate}
\end{frame}

\begin{frame}[t]{Teaching Plan}
\justifying
\begin{itemize}
\item First part: mathematical foundations with numerical examples.
\item Second part: complete derivation of the option-pricing model.
\item Final 30 minutes: guided exercises followed by solutions.
\end{itemize}\begin{alertblock}{Common mistake}Students should explain each symbol before substituting numbers.\end{alertblock}
\end{frame}

\section{9.1 Exponential Growth and the Number $e$}
\begin{frame}[plain]
\centering\vfill
{\color{TitleBlue}\Large\bfseries 9.1 Exponential Growth and the Number $e$}\par
\vspace{0.5cm}{\large Why continuous time naturally creates an exponential function}\par
\vfill\end{frame}

\begin{frame}[t]{Simple Interest}
\justifying
\begin{itemize}
\item Simple interest applies the rate only to the original principal.
\item For principal $P$, annual rate $r$, and time $t$, the future value is
\item \[FV=P(1+rt).\]
\item The interest earned during one period does not itself earn interest.
\end{itemize}
\end{frame}

\begin{frame}[t]{Compound Interest}
\justifying
\begin{itemize}
\item Compounding reinvests interest after every compounding date.
\item If the year contains $m$ compounding periods,
\item \[FV_m=P\left(1+\frac{r}{m}\right)^{mt}.\]
\item Increasing $m$ makes compounding more frequent.
\end{itemize}
\end{frame}

\begin{frame}[t]{The Limit that Defines $e$}
\justifying
Set $r=1$, $t=1$, and let the number of compounding periods increase:
\[\left(1+\frac{1}{m}\right)^m.\]
\begin{center}
\begin{tabular}{rrrrrr}\toprule
$m$&1&2&4&12&365\\\midrule
Value&2.0000&2.2500&2.4414&2.6130&2.7146\\\bottomrule
\end{tabular}
\end{center}
The sequence approaches
\[e=\lim_{m\to\infty}\left(1+\frac{1}{m}\right)^m\approx2.71828.\]
\end{frame}

\begin{frame}[t]{Why $e$ Appears in Finance}
\justifying
\begin{itemize}
\item Continuous compounding takes the limit as the compounding interval approaches zero.
\item \[\lim_{m\to\infty}P\left(1+\frac{r}{m}\right)^{mt}=Pe^{rt}.\]
\item The exponential function is the unique growth function whose instantaneous growth rate is proportional to its current level.
\item \[\frac{d}{dt}e^{rt}=re^{rt}.\]
\end{itemize}
\end{frame}

\begin{frame}[t]{Exponential Growth as a Differential Equation}
\justifying
Suppose value grows at instantaneous rate $r$:
\[\frac{dV_t}{dt}=rV_t.\]
Separate variables:
\[\frac{dV_t}{V_t}=r\,dt.\]
Integrate from $0$ to $t$:
\[\ln V_t-\ln V_0=rt.\]
Exponentiating gives
\[V_t=V_0e^{rt}.\]
\end{frame}

\begin{frame}[t]{Discounting with Continuous Compounding}
\justifying
If one dollar grows to $e^{rT}$ dollars at time $T$, then one dollar received at time $T$ has present value
\[PV=e^{-rT}.\]
For a certain payment $K$ at time $T$,
\[PV(K)=Ke^{-rT}.\]
\begin{block}{Finance interpretation}The term $e^{-rT}$ appears in Black--Scholes because the strike payment occurs at expiration and must be discounted continuously.\end{block}
\end{frame}

\begin{frame}[t]{Logarithms Undo Exponentials}
\justifying
\begin{itemize}
\item If $y=e^x$, then $x=\ln y$.
\item $\ln(ab)=\ln a+\ln b$ converts multiplication into addition.
\item $\ln(a^q)=q\ln a$ brings powers down as coefficients.
\item These properties make continuously compounded returns additive through time.
\end{itemize}\begin{block}{Example}If $S_T=S_0e^R$, then $R=\ln(S_T/S_0)$.\end{block}
\end{frame}

\section{9.2 Probability Foundations}
\begin{frame}[plain]
\centering\vfill
{\color{TitleBlue}\Large\bfseries 9.2 Probability Foundations}\par
\vspace{0.5cm}{\large The minimum probability language needed for stochastic calculus}\par
\vfill\end{frame}

\begin{frame}[t]{Random Variables}
\justifying
\begin{itemize}
\item A random variable assigns a number to each possible outcome.
\item A discrete random variable takes countable values.
\item A continuous random variable can take any value in an interval.
\item A stock return is uncertain today and becomes an observed number later.
\end{itemize}
\end{frame}

\begin{frame}[t]{Expected Value}
\justifying
For discrete outcomes,
\[\E[X]=\sum_i p_i x_i.\]
For a continuous variable with density $f(x)$,
\[\E[X]=\int_{-\infty}^{\infty}x f(x)\,dx.\]
Expected value is a probability-weighted average, not a guaranteed outcome.
\end{frame}

\begin{frame}[t]{Variance and Standard Deviation}
\justifying
Variance measures average squared deviation from the mean:
\[\Var(X)=\E[(X-\E[X])^2].\]
Standard deviation uses the original unit:
\[\sigma_X=\sqrt{\Var(X)}.\]
In an option model, volatility is the annualized standard deviation of returns.
\end{frame}

\begin{frame}[t]{Conditional Expectation}
\justifying
\begin{itemize}
\item $\E[X\mid\mathcal F_t]$ means the expected value of $X$ using information available at time $t$.
\item The information set grows as prices and news arrive.
\item Option value at time $t$ depends on the distribution of future payoffs conditional on current information.
\end{itemize}\begin{block}{Notation}$\mathcal F_t$ represents everything known by time $t$.\end{block}
\end{frame}

\begin{frame}[t]{Independence}
\justifying
\begin{itemize}
\item Two increments are independent when observing one provides no information about the other.
\item For independent $X$ and $Y$, $\Cov(X,Y)=0$.
\item Brownian motion assumes independent increments over nonoverlapping time intervals.
\end{itemize}\begin{alertblock}{Common mistake}Zero covariance does not always imply independence. Joint normality makes the implication valid.\end{alertblock}
\end{frame}

\begin{frame}[t]{The Central Limit Idea}
\justifying
\begin{itemize}
\item A long-period return combines many small shocks.
\item Under suitable conditions, a standardized sum of independent small shocks approaches a normal distribution.
\item This limiting idea motivates normal increments in Brownian motion.
\item Market returns are not exactly normal, especially in crises.
\end{itemize}
\end{frame}

\section{9.3 Normal and Lognormal Distributions}
\begin{frame}[plain]
\centering\vfill
{\color{TitleBlue}\Large\bfseries 9.3 Normal and Lognormal Distributions}\par
\vspace{0.5cm}{\large Why the model applies normality to log returns rather than prices}\par
\vfill\end{frame}

\begin{frame}[t]{The Standard Normal Distribution}
\justifying
A standard normal variable $Z$ has density
\[\phi(z)=\frac{1}{\sqrt{2\pi}}e^{-z^2/2}.\]
Its cumulative distribution function is
\[N(z)=\Pr(Z\le z)=\int_{-\infty}^{z}\phi(u)\,du.\]
The function $N(\cdot)$ appears in the Black--Scholes formula.
\end{frame}

\begin{frame}[t]{Why the Normal Density Contains $e$}
\justifying
A bell-shaped density with derivative proportional to $-z$ times itself satisfies
\[\phi'(z)=-z\phi(z).\]
Solving this differential equation gives
\[\phi(z)=Ce^{-z^2/2}.\]
The constant $C=1/\sqrt{2\pi}$ makes the total area equal to one.
\end{frame}

\begin{frame}[t]{Normal Variables Can Be Negative}
\justifying
\begin{itemize}
\item A normal variable ranges from $-\infty$ to $+\infty$.
\item Modeling the stock price itself as normal assigns positive probability to a negative price.
\item Limited liability prevents an ordinary share price from falling below zero.
\item The model therefore places a normal distribution on the logarithm of price.
\end{itemize}
\end{frame}

\begin{frame}[t]{Definition of a Lognormal Variable}
\justifying
If
\[Y=\ln X\sim N(m,v),\]
then
\[X=e^Y\]
has a lognormal distribution. Because the exponential function is positive,
\[X>0.\]
The distribution is right-skewed and allows large upside outcomes.
\end{frame}

\begin{frame}[t]{Mean of a Lognormal Variable}
\justifying
If $Y\sim N(m,v)$ and $X=e^Y$, then
\[\E[X]=e^{m+v/2}.\]
The term $v/2$ matters. In the stock model, this term creates the $-\sigma^2/2$ adjustment in expected log return.
\end{frame}

\begin{frame}[t]{Price Ratios and Log Returns}
\justifying
The continuously compounded return is
\[R_{t,T}=\ln\left(\frac{S_T}{S_t}\right).\]
If log return is normal, then
\[\frac{S_T}{S_t}=e^{R_{t,T}}\]
is lognormal and remains positive.
\end{frame}

\begin{frame}[t]{Why Log Returns Are Convenient}
\justifying
\begin{itemize}
\item Additivity: $\ln(S_2/S_0)=\ln(S_1/S_0)+\ln(S_2/S_1)$.
\item Positivity: exponentiating a real number gives a positive price.
\item Scale: percentage changes matter more than equal dollar changes across different price levels.
\item Aggregation: sums of many small log-return shocks support a normal approximation.
\end{itemize}
\end{frame}

\begin{frame}[t]{Modeling Choice, Not a Natural Law}
\justifying
\begin{itemize}
\item Actual returns show jumps, volatility clustering, skewness, and heavy tails.
\item The lognormal assumption creates a tractable benchmark.
\item Practitioners often use implied volatility, stochastic volatility, or jump models when the benchmark is inadequate.
\end{itemize}\begin{alertblock}{Common mistake}The model does not claim that every observed return is normally distributed.\end{alertblock}
\end{frame}

\section{9.4 Random Walks and Brownian Motion}
\begin{frame}[plain]
\centering\vfill
{\color{TitleBlue}\Large\bfseries 9.4 Random Walks and Brownian Motion}\par
\vspace{0.5cm}{\large From many small shocks to a continuous random path}\par
\vfill\end{frame}

\begin{frame}[t]{A Discrete Random Walk}
\justifying
Let each one-period shock be $\varepsilon_i\in\{-1,+1\}$ with equal probability. Define
\[X_n=\sum_{i=1}^{n}\varepsilon_i.\]
Then
\[\E[X_n]=0,\qquad \Var(X_n)=n.\]
The uncertainty grows with the number of steps.
\end{frame}

\begin{frame}[t]{Scaling a Random Walk}
\justifying
Divide a fixed horizon $T$ into $n$ intervals of length
\[\Delta t=T/n.\]
Scale each shock by $\sqrt{\Delta t}$:
\[\Delta W_i=\sqrt{\Delta t}\,\varepsilon_i.\]
Then
\[\Var\left(\sum_{i=1}^n\Delta W_i\right)=n\Delta t=T.\]
\end{frame}

\begin{frame}[t]{Why the Shock Scales with $\sqrt{dt}$}
\justifying
If each shock scaled with $dt$, total variance would be
\[n(dt)^2=T\,dt\to0.\]
All randomness would disappear. If each shock had fixed variance, total variance would explode. Scaling by $\sqrt{dt}$ gives
\[n(\sqrt{dt})^2=n\,dt=T,\]
a finite nonzero variance.
\end{frame}

\begin{frame}[t]{Brownian Motion}
\justifying
\begin{itemize}
\item $W_0=0$.
\item Increments over nonoverlapping intervals are independent.
\item $W_{t+h}-W_t\sim N(0,h)$.
\item Sample paths are continuous with probability one.
\item The path is continuous but almost surely nowhere differentiable.
\end{itemize}
\end{frame}

\begin{frame}[t]{Stationary Increments}
\justifying
For any $t$ and $h>0$,
\[W_{t+h}-W_t\sim N(0,h).\]
Only the interval length $h$ matters. The calendar starting point $t$ does not change the increment distribution.
\[\E[W_{t+h}-W_t]=0,\qquad \Var(W_{t+h}-W_t)=h.\]
\end{frame}

\begin{frame}[t]{Why Brownian Motion Appears in Finance}
\justifying
\begin{itemize}
\item Price changes reflect many small pieces of information.
\item The aggregate of many small shocks suggests a continuous Gaussian limit.
\item Independent increments create a mathematically manageable benchmark.
\item The model provides continuous trading and closed-form valuation.
\end{itemize}\begin{alertblock}{Common mistake}This explanation motivates the assumption. It does not prove that real prices must follow Brownian motion.\end{alertblock}
\end{frame}

\begin{frame}[t]{Brownian Paths Are Rough}
\justifying
\begin{itemize}
\item Over a short interval, the typical increment has size $\sqrt{dt}$.
\item The ratio $dW/dt$ has typical size $1/\sqrt{dt}$ and becomes unstable as $dt\to0$.
\item Ordinary derivatives do not exist along a Brownian path.
\item Stochastic calculus defines integrals directly rather than differentiating the path.
\end{itemize}
\end{frame}

\begin{frame}[t]{Quadratic Variation}
\justifying
For a partition $0=t_0<\cdots<t_n=T$,
\[\sum_{i=1}^{n}(W_{t_i}-W_{t_{i-1}})^2\to T.\]
This nonzero limit is Brownian quadratic variation. In differential shorthand,
\[(dW)^2=dt.\]
This identity creates the second-derivative term in It\^o\textquotesingle s lemma.
\end{frame}

\section{9.5 Stochastic Differential Notation}
\begin{frame}[plain]
\centering\vfill
{\color{TitleBlue}\Large\bfseries 9.5 Stochastic Differential Notation}\par
\vspace{0.5cm}{\large What $dt$, $dW$, and $dS$ actually mean}\par
\vfill\end{frame}

\begin{frame}[t]{Meaning of $dt$}
\justifying
\begin{itemize}
\item $dt$ denotes a very small time interval in differential notation.
\item In a rigorous construction, equations describe limits of sums as the largest interval approaches zero.
\item A deterministic change proportional to $dt$ accumulates linearly with time.
\end{itemize}
\end{frame}

\begin{frame}[t]{Meaning of $dW$}
\justifying
Over a short interval $dt$,
\[dW\sim N(0,dt).\]
Equivalently,
\[dW=\sqrt{dt}\,Z,\qquad Z\sim N(0,1).\]
Therefore $dW$ has order $\sqrt{dt}$, which is larger than $dt$ over a very short interval.
\end{frame}

\begin{frame}[t]{The It\^o Multiplication Rules}
\justifying
The orders of small quantities imply
\[(dt)^2=0,\qquad dt\,dW=0,\qquad (dW)^2=dt.\]
The first two terms vanish relative to $dt$. The last term remains because Brownian motion has nonzero quadratic variation.
\end{frame}

\begin{frame}[t]{Meaning of a Stochastic Differential Equation}
\justifying
The notation
\[dX_t=a(X_t,t)dt+b(X_t,t)dW_t\]
is shorthand for the integral equation
\[X_t=X_0+\int_0^t a(X_s,s)ds+\int_0^t b(X_s,s)dW_s.\]
The first integral is ordinary. The second is an It\^o stochastic integral.
\end{frame}

\begin{frame}[t]{Drift and Diffusion}
\justifying
\begin{itemize}
\item $a(X_t,t)dt$ is the conditional expected change over a short interval.
\item $b(X_t,t)dW_t$ is the unpredictable shock.
\item The drift has order $dt$. The diffusion has order $\sqrt{dt}$.
\item At very high frequency, random variation dominates the local path.
\end{itemize}
\end{frame}

\begin{frame}[t]{Euler Discretization}
\justifying
For a small finite step $\Delta t$,
\[X_{t+\Delta t}\approx X_t+a(X_t,t)\Delta t+b(X_t,t)\sqrt{\Delta t}\,Z.\]
This approximation connects stochastic notation to simulation and spreadsheets.
\end{frame}

\section{9.6 Geometric Brownian Motion}
\begin{frame}[plain]
\centering\vfill
{\color{TitleBlue}\Large\bfseries 9.6 Geometric Brownian Motion}\par
\vspace{0.5cm}{\large The stock-price process used in Black--Scholes}\par
\vfill\end{frame}

\begin{frame}[t]{The Stock-Price SDE}
\justifying
Black--Scholes assumes
\[dS_t=\mu S_tdt+\sigma S_tdW_t.\]
Dividing by $S_t$ gives
\[\frac{dS_t}{S_t}=\mu dt+\sigma dW_t.\]
$\mu$ is the instantaneous expected return and $\sigma$ is volatility.
\end{frame}

\begin{frame}[t]{Why the Model Is Multiplicative}
\justifying
\begin{itemize}
\item The dollar volatility $\sigma S_t$ rises with the price level.
\item The model describes proportional rather than fixed dollar changes.
\item The same percentage shock has a larger dollar effect on a more valuable share.
\item The resulting price remains positive.
\end{itemize}
\end{frame}

\begin{frame}[t]{A Tempting but Incorrect Log Step}
\justifying
One might incorrectly write
\[d\ln S_t=\frac{dS_t}{S_t}.\]
This ordinary-calculus rule misses Brownian quadratic variation. The correct expression contains an additional term:
\[d\ln S_t=\left(\mu-\frac{1}{2}\sigma^2\right)dt+\sigma dW_t.\]
\end{frame}

\begin{frame}[t]{Solution of Geometric Brownian Motion}
\justifying
Integrate the correct log equation from $0$ to $T$:
\[\ln\frac{S_T}{S_0}=\left(\mu-\frac{1}{2}\sigma^2\right)T+\sigma W_T.\]
Exponentiating,
\[S_T=S_0\exp\left[\left(\mu-\frac{1}{2}\sigma^2\right)T+\sigma W_T\right].\]
\end{frame}

\begin{frame}[t]{Distribution of the Future Price}
\justifying
Because $W_T\sim N(0,T)$,
\[\ln S_T\sim N\left(\ln S_0+\left(\mu-\frac{1}{2}\sigma^2\right)T,\ \sigma^2T\right).\]
Therefore $S_T$ is lognormal. Its expectation is
\[\E[S_T]=S_0e^{\mu T}.\]
\end{frame}

\section{9.7 It\^o\textquotesingle s Lemma}
\begin{frame}[plain]
\centering\vfill
{\color{TitleBlue}\Large\bfseries 9.7 It\^o\textquotesingle s Lemma}\par
\vspace{0.5cm}{\large The stochastic version of the chain rule}\par
\vfill\end{frame}

\begin{frame}[t]{Ordinary Taylor Expansion}
\justifying
For a smooth function $f(x,t)$,
\[df\approx f_tdt+f_xdx+\frac{1}{2}f_{xx}(dx)^2+f_{xt}dxdt+\frac{1}{2}f_{tt}(dt)^2+\cdots\]
In ordinary calculus, $(dx)^2$ usually vanishes faster than $dt$. Brownian motion changes this ordering.
\end{frame}

\begin{frame}[t]{Insert a Stochastic Change}
\justifying
Let
\[dX=a\,dt+b\,dW.\]
Then
\[(dX)^2=a^2(dt)^2+2ab\,dt\,dW+b^2(dW)^2.\]
Using the It\^o rules,
\[(dX)^2=b^2dt.\]
\end{frame}

\begin{frame}[t]{Terms That Survive}
\justifying
In the Taylor expansion:
\[(dt)^2\text{ has order }dt^2\to0,\]
\[dt\,dW\text{ has order }dt^{3/2}\to0,\]
\[(dW)^2\text{ has order }dt\text{ and remains.}\]
Therefore the second derivative with respect to $x$ cannot be discarded.
\end{frame}

\begin{frame}[t]{It\^o\textquotesingle s Lemma in One Dimension}
\justifying
If
\[dX=a(X,t)dt+b(X,t)dW\]
and $f=f(X,t)$, then
\[df=\left(f_t+af_x+\frac{1}{2}b^2f_{xx}\right)dt+bf_xdW.\]\begin{block}{Key difference from ordinary calculus}The correction term is $\tfrac12 b^2 f_{xx}dt$.\end{block}
\end{frame}

\begin{frame}[t]{Why the Factor $1/2$ Appears}
\justifying
\begin{itemize}
\item The coefficient comes directly from the second-order Taylor term.
\item Taylor expansion contains $\tfrac12 f_{xx}(dX)^2$.
\item Brownian quadratic variation converts $(dX)^2$ into $b^2dt$.
\item The surviving term becomes $\tfrac12 b^2f_{xx}dt$.
\end{itemize}
\end{frame}

\begin{frame}[t]{Example: $f(X)=X^2$}
\justifying
Here $f_x=2X$ and $f_{xx}=2$. If $dX=a,dt+b,dW$, then
\[d(X^2)=2X\,dX+b^2dt.\]
Substituting $dX$,
\[d(X^2)=(2aX+b^2)dt+2bX\,dW.\]
Ordinary calculus would miss the $b^2dt$ term.
\end{frame}

\begin{frame}[t]{Example: $f(X)=\ln X$}
\justifying
For $f(X)=\ln X$,
\[f_x=\frac1X,\qquad f_{xx}=-\frac1{X^2}.\]
With $dS=\mu Sdt+\sigma SdW$,
\[d\ln S=\frac1S dS-\frac12\frac1{S^2}(\sigma^2S^2)dt.\]
Therefore
\[d\ln S=\left(\mu-\frac12\sigma^2\right)dt+\sigma dW.\]
\end{frame}

\begin{frame}[t]{Example: $f(X)=e^X$}
\justifying
For $f(X)=e^X$, both first and second derivatives equal $e^X$. Thus
\[d(e^X)=e^X\left(a+\frac12b^2\right)dt+e^Xb\,dW.\]
The convexity of the exponential function raises the drift by $b^2/2$.
\end{frame}

\begin{frame}[t]{It\^o\textquotesingle s Lemma for an Option}
\justifying
Let the option value be $V(S,t)$ and let
\[dS=\mu Sdt+\sigma SdW.\]
Then
\[dV=\left(V_t+\mu SV_S+\frac12\sigma^2S^2V_{SS}\right)dt+\sigma SV_SdW.\]
This equation is the bridge from the stock process to option pricing.
\end{frame}

\begin{frame}[t]{Meaning of the Derivatives}
\justifying
\begin{itemize}
\item $V_t$ measures time decay with the stock price fixed.
\item $V_S=\Delta$ measures first-order exposure to the stock price.
\item $V_{SS}=\Gamma$ measures curvature of the option value.
\item The term $\tfrac12\sigma^2S^2\Gamma$ is the It\^o correction caused by price uncertainty.
\end{itemize}
\end{frame}

\section{9.8 No-Arbitrage and Replication}
\begin{frame}[plain]
\centering\vfill
{\color{TitleBlue}\Large\bfseries 9.8 No-Arbitrage and Replication}\par
\vspace{0.5cm}{\large Why expected stock return disappears from option value}\par
\vfill\end{frame}

\begin{frame}[t]{Law of One Price}
\justifying
\begin{itemize}
\item Two portfolios with identical future cash flows must have the same current value.
\item If prices differ, traders buy the cheaper portfolio and sell the more expensive one.
\item A persistent price difference would create an arbitrage profit without net risk or investment.
\end{itemize}
\end{frame}

\begin{frame}[t]{A Locally Riskless Portfolio}
\justifying
Hold $\Delta$ shares and short one option:
\[\Pi=\Delta S-V.\]
Over a short interval,
\[d\Pi=\Delta dS-dV.\]
Choose $\Delta$ so that the $dW$ exposure cancels.
\end{frame}

\begin{frame}[t]{Delta Hedge}
\justifying
The random term in $dV$ is
\[\sigma SV_SdW.\]
The random term in $\Delta dS$ is
\[\Delta\sigma SdW.\]
Set
\[\Delta=V_S.\]
The two random terms cancel locally.
\end{frame}

\begin{frame}[t]{Return on the Hedged Portfolio}
\justifying
After choosing $\Delta=V_S$,
\[d\Pi=\left(-V_t-\frac12\sigma^2S^2V_{SS}\right)dt.\]
The portfolio has no $dW$ term over the instant. No arbitrage therefore requires
\[d\Pi=r\Pi dt.\]
\end{frame}

\begin{frame}[t]{Why $\mu$ Disappears}
\justifying
\begin{itemize}
\item Both the stock and option contain exposure to the same Brownian shock.
\item The delta hedge removes that shock and also removes the stock drift term from the portfolio equation.
\item The riskless portfolio earns the risk-free rate.
\item Option value therefore depends on $r$ rather than the stock’s expected return $\mu$.
\end{itemize}
\end{frame}

\begin{frame}[t]{Dynamic Replication}
\justifying
\begin{itemize}
\item Option delta changes when the stock price or time changes.
\item The replicating portfolio must rebalance continuously in the theoretical model.
\item Real trading is discrete and costly, so actual hedges retain residual risk.
\item Gamma measures how quickly delta changes with the stock price.
\end{itemize}
\end{frame}

\section{9.9 The Black--Scholes Partial Differential Equation}
\begin{frame}[plain]
\centering\vfill
{\color{TitleBlue}\Large\bfseries 9.9 The Black--Scholes Partial Differential Equation}\par
\vspace{0.5cm}{\large Complete algebra from the hedge to the PDE}\par
\vfill\end{frame}

\begin{frame}[t]{Step 1: Stock Dynamics}
\justifying
Start with
\[dS=\mu Sdt+\sigma SdW.\]
The stock has expected proportional change $\mu dt$ and random proportional change $\sigma dW$.
\end{frame}

\begin{frame}[t]{Step 2: Option Dynamics}
\justifying
It\^o\textquotesingle s lemma gives
\[dV=\left(V_t+\mu SV_S+\frac12\sigma^2S^2V_{SS}\right)dt+\sigma SV_SdW.\]
\end{frame}

\begin{frame}[t]{Step 3: Form the Portfolio}
\justifying
Set
\[\Pi=V_S S-V.\]
Over an instant, hold $V_S$ fixed while computing the portfolio change:
\[d\Pi=V_SdS-dV.\]
\end{frame}

\begin{frame}[t]{Step 4: Cancel the Random Shock}
\justifying
\small
Substitute $dS$ and $dV$:
\[d\Pi=V_S(\mu Sdt+\sigma SdW)-\left[\left(V_t+\mu SV_S+\frac12\sigma^2S^2V_{SS}\right)dt+\sigma SV_SdW\right].\]
The $\mu SV_Sdt$ and $\sigma SV_SdW$ terms cancel.
\end{frame}

\begin{frame}[t]{Step 5: Riskless Portfolio Change}
\justifying
After cancellation,
\[d\Pi=-\left(V_t+\frac12\sigma^2S^2V_{SS}\right)dt.\]
The portfolio value is
\[\Pi=SV_S-V.\]
\end{frame}

\begin{frame}[t]{Step 6: Apply No Arbitrage}
\justifying
A locally riskless self-financing portfolio earns the risk-free rate:
\[-\left(V_t+\frac12\sigma^2S^2V_{SS}\right)dt=r(SV_S-V)dt.\]
Move every term to the left:
\[\boxed{V_t+\frac12\sigma^2S^2V_{SS}+rSV_S-rV=0}.\]
\end{frame}

\begin{frame}[t]{Boundary and Terminal Conditions}
\justifying
The PDE alone does not identify a unique option price. A European call requires
\[V(S,T)=\max(S-K,0).\]
Useful boundaries are
\[V(0,t)=0,\qquad V(S,t)\sim S-Ke^{-r(T-t)}\text{ as }S\to\infty.\]
\end{frame}

\begin{frame}[t]{Economic Meaning of the PDE}
\justifying
\begin{itemize}
\item $V_t$ captures the passage of time.
\item $\tfrac12\sigma^2S^2V_{SS}$ prices curvature under uncertainty.
\item $rSV_S$ reflects the financing of the stock component.
\item $-rV$ reflects the required return on the option value after hedging.
\end{itemize}
\end{frame}

\section{9.10 Solving the Black--Scholes PDE}
\begin{frame}[plain]
\centering\vfill
{\color{TitleBlue}\Large\bfseries 9.10 Solving the Black--Scholes PDE}\par
\vspace{0.5cm}{\large From the pricing equation to the normal distribution}\par
\vfill\end{frame}

\begin{frame}[t]{Change of Variables}
\justifying
Define time to maturity and log moneyness:
\[\tau=\frac12\sigma^2(T-t),\qquad x=\ln(S/K).\]
Scale the option value:
\[v(x,\tau)=V(S,t)/K.\]
The logarithm converts the multiplicative stock variable into an additive variable.
\end{frame}

\begin{frame}[t]{Derivative Transformation}
\justifying
Because $x=\ln(S/K)$,
\[SV_S=Kv_x,\qquad S^2V_{SS}=K(v_{xx}-v_x).\]
Because $\tau=\tfrac12\sigma^2(T-t)$,
\[V_t=-\frac12\sigma^2K v_\tau.\]
\end{frame}

\begin{frame}[t]{Dimensionless PDE}
\justifying
Let
\[k=\frac{2r}{\sigma^2}.\]
Substitution into the Black--Scholes PDE gives
\[v_\tau=v_{xx}+(k-1)v_x-kv.\]
This equation still contains a first derivative and a level term.
\end{frame}

\begin{frame}[t]{Remove the Extra Terms}
\justifying
Set
\[v(x,\tau)=e^{ax+b\tau}u(x,\tau).\]
Choose
\[a=\frac{1-k}{2},\qquad b=-\frac{(k+1)^2}{4}.\]
These constants cancel the $u_x$ and $u$ terms.
\end{frame}

\begin{frame}[t]{The Heat Equation}
\justifying
After the transformation,
\[u_\tau=u_{xx}.\]
The Black--Scholes pricing problem has become the classical heat equation. Diffusion of heat and diffusion of probability share the same mathematical structure.
\end{frame}

\begin{frame}[t]{Heat Kernel}
\justifying
\small
The heat equation solution averages the initial condition with a Gaussian kernel:
\[u(x,\tau)=\frac{1}{\sqrt{4\pi\tau}}\int_{-\infty}^{\infty}u(y,0)\exp\left[-\frac{(x-y)^2}{4\tau}\right]dy.\]
This Gaussian integral explains why the normal CDF enters the final option formula.
\end{frame}

\begin{frame}[t]{Terminal Payoff under the Transformation}
\justifying
At maturity, $\tau=0$ and
\[v(x,0)=\max(e^x-1,0).\]
The payoff is positive only when $x>0$, which means $S>K$. The heat-kernel integral therefore splits at the exercise boundary.
\end{frame}

\begin{frame}[t]{An Equivalent Risk-Neutral Route}
\justifying
The PDE solution equals a discounted expectation under a risk-neutral probability measure:
\[V_t=e^{-r(T-t)}\E^{\mathbb Q}[\text{payoff}\mid\mathcal F_t].\]
Under $\mathbb Q$,
\[dS=rSdt+\sigma SdW^{\mathbb Q}.\]
The stock drift changes from $\mu$ to $r$ for pricing.
\end{frame}

\section{9.11 The Black--Scholes Formula}
\begin{frame}[plain]
\centering\vfill
{\color{TitleBlue}\Large\bfseries 9.11 The Black--Scholes Formula}\par
\vspace{0.5cm}{\large Evaluation of the discounted payoff integral}\par
\vfill\end{frame}

\begin{frame}[t]{Risk-Neutral Distribution of $S_T$}
\justifying
Let $\tau=T-t$. Under the risk-neutral measure,
\[\ln S_T\sim N\left(\ln S_t+\left(r-\frac12\sigma^2\right)\tau,\ \sigma^2\tau\right).\]
Equivalently,
\[S_T=S_t\exp\left[\left(r-\frac12\sigma^2\right)\tau+\sigma\sqrt\tau Z\right].\]
\end{frame}

\begin{frame}[t]{Split the Call Payoff}
\justifying
A European call payoff is
\[(S_T-K)^+=S_T\mathbf 1_{\{S_T>K\}}-K\mathbf 1_{\{S_T>K\}}.\]
Therefore
\[C_t=e^{-r\tau}\E^{\mathbb Q}[S_T\mathbf 1_{\{S_T>K\}}]-Ke^{-r\tau}\Pr^{\mathbb Q}(S_T>K).\]
\end{frame}

\begin{frame}[t]{Exercise Probability and $d_2$}
\justifying
The event $S_T>K$ becomes
\[Z>\frac{\ln(K/S_t)-(r-\sigma^2/2)\tau}{\sigma\sqrt\tau}.\]
Using normal symmetry,
\[\Pr^{\mathbb Q}(S_T>K)=N(d_2),\]
where
\[d_2=\frac{\ln(S_t/K)+(r-\sigma^2/2)\tau}{\sigma\sqrt\tau}.\]
\end{frame}

\begin{frame}[t]{Stock-Weighted Term and $d_1$}
\justifying
Evaluating the truncated lognormal expectation gives
\[e^{-r\tau}\E^{\mathbb Q}[S_T\mathbf 1_{\{S_T>K\}}]=S_tN(d_1),\]
where
\[d_1=\frac{\ln(S_t/K)+(r+\sigma^2/2)\tau}{\sigma\sqrt\tau}=d_2+\sigma\sqrt\tau.\]
\end{frame}

\begin{frame}[t]{European Call Formula}
\justifying
Combining the two payoff components gives
\[\boxed{C_t=S_tN(d_1)-Ke^{-r\tau}N(d_2)}.\]
The first term is the option-adjusted value of receiving the stock. The second term is the option-adjusted present value of paying the strike.
\end{frame}

\begin{frame}[t]{European Put Formula}
\justifying
Put--call parity gives
\[C_t-P_t=S_t-Ke^{-r\tau}.\]
Therefore
\[\boxed{P_t=Ke^{-r\tau}N(-d_2)-S_tN(-d_1)}.\]
\end{frame}

\begin{frame}[t]{Interpretation of $N(d_1)$ and $N(d_2)$}
\justifying
\begin{itemize}
\item $N(d_1)$ is the call delta for a non-dividend-paying stock.
\item $N(d_2)$ is the risk-neutral probability that the call finishes in the money.
\item The two numbers differ because the first expectation weights outcomes by the future stock price.
\item $d_1=d_2+\sigma\sqrt\tau$.
\end{itemize}
\end{frame}

\begin{frame}[t]{Why the Formula Contains Two Uses of $e$}
\justifying
\begin{itemize}
\item $e^{-r\tau}$ discounts the future strike payment continuously.
\item The lognormal stock model comes from exponentiating a normal log return.
\item Both appearances trace back to continuous proportional growth.
\item The normal density itself also contains an exponential term.
\end{itemize}
\end{frame}

\section{9.12 Interpretation, Extensions, and Limits}
\begin{frame}[plain]
\centering\vfill
{\color{TitleBlue}\Large\bfseries 9.12 Interpretation, Extensions, and Limits}\par
\vspace{0.5cm}{\large What the model assumes and how practitioners adapt it}\par
\vfill\end{frame}

\begin{frame}[t]{Black--Scholes Assumptions}
\justifying
\begin{itemize}
\item The stock follows geometric Brownian motion with constant volatility.
\item Markets permit continuous trading without taxes or transaction costs.
\item Investors can borrow and lend at the risk-free rate.
\item Short selling is possible and securities are divisible.
\item The option is European and the basic stock pays no dividend.
\end{itemize}
\end{frame}

\begin{frame}[t]{Dividend Yield Extension}
\justifying
For a continuous dividend yield $q$, risk-neutral stock dynamics become
\[dS=(r-q)Sdt+\sigma SdW^{\mathbb Q}.\]
The call price becomes
\[C=S e^{-q\tau}N(d_1)-Ke^{-r\tau}N(d_2),\]
with
\[d_1=\frac{\ln(S/K)+(r-q+\sigma^2/2)\tau}{\sigma\sqrt\tau}.\]
\end{frame}

\begin{frame}[t]{The Greeks}
\justifying
\begin{center}
\begin{tabular}{lll}\toprule
Greek&Derivative&Main exposure\\\midrule
Delta&$V_S$&stock-price sensitivity\\
Gamma&$V_{SS}$&change in delta\\
Theta&$V_t$&passage of time\\
Vega&$V_\sigma$&volatility sensitivity\\
Rho&$V_r$&interest-rate sensitivity\\\bottomrule
\end{tabular}
\end{center}
\end{frame}

\begin{frame}[t]{Volatility and Option Value}
\justifying
\begin{itemize}
\item A call holder has limited downside and retains upside.
\item Greater volatility spreads the distribution of future prices.
\item The favorable tail becomes more valuable while the payoff floor remains zero.
\item Both call and put values generally rise with volatility.
\end{itemize}
\end{frame}

\begin{frame}[t]{Implied Volatility}
\justifying
Market price can be inserted into the Black--Scholes formula and solved for $\sigma$:
\[C_{market}=BS(S,K,r,\tau,\sigma_{impl}).\]
Implied volatility summarizes the volatility level consistent with the observed option price.
\end{frame}

\begin{frame}[t]{Volatility Smile}
\justifying
\begin{itemize}
\item Constant volatility predicts the same volatility for all strikes and maturities.
\item Observed implied volatility commonly varies across strike and maturity.
\item The pattern reflects skewness, heavy tails, jumps, supply and demand, and market risk premiums.
\item The smile shows where the benchmark model misses market behavior.
\end{itemize}
\end{frame}

\begin{frame}[t]{Global Financial Strategy Applications}
\justifying
\begin{itemize}
\item Currency options protect uncertain foreign-currency cash flows while preserving favorable exchange-rate movements.
\item Equity can be viewed as an option on firm assets when debt is outstanding.
\item Investment timing creates real options when management can wait for country or regulatory information.
\item Option-based compensation affects dilution and managerial risk incentives.
\end{itemize}
\end{frame}

\begin{frame}[t]{Model Risk}
\justifying
\begin{itemize}
\item Volatility changes over time and cannot be observed directly.
\item Prices jump after major announcements or policy changes.
\item Continuous hedging is impossible in practice.
\item Liquidity, transaction costs, and funding constraints become important during stress.
\item The formula is a benchmark that requires judgment and stress testing.
\end{itemize}
\end{frame}

\begin{frame}[t]{Full Derivation Map}
\justifying
\begin{enumerate}
\item Continuous proportional growth produces $e$ and logarithms.
\item Many small shocks motivate Brownian motion and normal increments.
\item Exponentiating normal log returns produces positive lognormal prices.
\item Brownian quadratic variation produces It\^o\textquotesingle s correction term.
\item Delta hedging removes the random shock and the unknown stock drift.
\item No arbitrage produces the PDE.
\item Gaussian integration produces $N(d_1)$ and $N(d_2)$.
\end{enumerate}
\end{frame}

\section{9.P 30-Minute Practice}
\begin{frame}[plain]
\centering\vfill
{\color{TitleBlue}\Large\bfseries 9.P 30-Minute Practice}\par
\vspace{0.5cm}{\large Foundations, derivation, and interpretation}\par
\vfill\end{frame}

\begin{frame}[t]{Practice Instructions}
\justifying
\begin{itemize}
\item Work individually for 20 minutes.
\item Compare derivations with a partner for 5 minutes.
\item Use the final 5 minutes to discuss economic interpretation.
\item Show the order of each small term when applying It\^o calculus.
\end{itemize}
\end{frame}

\begin{frame}[t]{Practice 1: Continuous Compounding}
\justifying
A principal of \$1,000 earns a continuously compounded annual rate of 5\% for three years. Calculate the future value and explain why $e$ appears.
\end{frame}

\begin{frame}[t]{Practice 2: Log Return}
\justifying
A stock rises from \$80 to \$92. Compute the continuously compounded return and compare it with the simple return.
\end{frame}

\begin{frame}[t]{Practice 3: Brownian Increment}
\justifying
Let $dt=0.01$ year and $dW=\sqrt{dt}Z$. If $Z=-0.6$, compute $dW$. Compare the sizes of $dt$, $dW$, $(dW)^2$, and $dt\,dW$.
\end{frame}

\begin{frame}[t]{Practice 4: One It\^o Step}
\justifying
If $dX=2dt+3dW$, use It\^o\textquotesingle s lemma to find $d(X^2)$.
\end{frame}

\begin{frame}[t]{Practice 5: Log Stock Price}
\justifying
Starting from $dS=\mu Sdt+\sigma SdW$, derive $d\ln S$ and identify the It\^o correction.
\end{frame}

\begin{frame}[t]{Practice 6: Delta Hedge}
\justifying
An option has delta $0.65$. Construct a portfolio using shares and one short option that eliminates the instantaneous $dW$ exposure.
\end{frame}

\begin{frame}[t]{Practice 7: PDE Algebra}
\justifying
Starting with $\Pi=V_SS-V$, reproduce the algebra leading to the Black--Scholes PDE.
\end{frame}

\begin{frame}[t]{Practice 8: Compute $d_1$ and $d_2$}
\justifying
Let $S=100$, $K=100$, $r=4\%$, $\sigma=20\%$, and $\tau=1$. Compute $d_1$ and $d_2$.
\end{frame}

\begin{frame}[t]{Practice 9: European Call}
\justifying
Using $N(0.30)=0.6179$ and $N(0.10)=0.5398$, calculate the call value for the data in Practice 8.
\end{frame}

\begin{frame}[t]{Practice 10: Assumption Diagnosis}
\justifying
A currency suddenly jumps 12\% after a policy announcement. Identify the Black--Scholes assumptions challenged by this event and explain the likely hedging consequence.
\end{frame}

\section{9.S Solutions and Debrief}
\begin{frame}[plain]
\centering\vfill
{\color{TitleBlue}\Large\bfseries 9.S Solutions and Debrief}\par
\vspace{0.5cm}{\large Detailed answers}\par
\vfill\end{frame}

\begin{frame}[t]{Solution 1}
\justifying
\[FV=1{,}000e^{0.05(3)}=\boxed{1{,}161.83}.\]
$e$ appears because continuous compounding is the limit of increasingly frequent discrete compounding.
\end{frame}

\begin{frame}[t]{Solution 2}
\justifying
\[R_{log}=\ln(92/80)=\boxed{13.976\%}.\]
\[R_{simple}=92/80-1=\boxed{15\%}.\]
The relationship is $R_{log}=\ln(1+R_{simple})$.
\end{frame}

\begin{frame}[t]{Solution 3}
\justifying
\[dW=\sqrt{0.01}(-0.6)=\boxed{-0.06}.\]
Orders: $dt=0.01$, $dW=-0.06$, $(dW)^2=0.0036$ for this realization, and $dt\,dW=-0.0006$. In stochastic-order reasoning, $(dW)^2$ accumulates like $dt$, while $dt\,dW$ vanishes.
\end{frame}

\begin{frame}[t]{Solution 4}
\justifying
For $f(X)=X^2$, $f_X=2X$ and $f_{XX}=2$:
\[d(X^2)=2X(2dt+3dW)+\frac12(2)(3^2)dt.\]
\[\boxed{d(X^2)=(4X+9)dt+6X dW}.\]
\end{frame}

\begin{frame}[t]{Solution 5}
\justifying
For $f(S)=\ln S$, $f_S=1/S$ and $f_{SS}=-1/S^2$:
\[d\ln S=\frac1S dS+\frac12\left(-\frac1{S^2}\right)(\sigma^2S^2)dt.\]
\[\boxed{d\ln S=(\mu-\sigma^2/2)dt+\sigma dW}.\]
\end{frame}

\begin{frame}[t]{Solution 6}
\justifying
Short one option and buy $0.65$ shares:
\[\Pi=0.65S-V.\]
\begin{itemize}
\item The stock contributes $0.65\sigma S dW$.
\item The short option contributes $-V_S\sigma S dW=-0.65\sigma S dW$.
\item The random terms cancel locally.
\end{itemize}
\end{frame}

\begin{frame}[t]{Solution 7}
\justifying
\[d\Pi=V_SdS-dV=-\left(V_t+\frac12\sigma^2S^2V_{SS}\right)dt.\]
No arbitrage requires $d\Pi=r(V_SS-V)dt$. Rearranging,
\[\boxed{V_t+\frac12\sigma^2S^2V_{SS}+rSV_S-rV=0}.\]
\end{frame}

\begin{frame}[t]{Solution 8}
\justifying
\[d_1=\frac{\ln(1)+(0.04+0.20^2/2)(1)}{0.20}=\boxed{0.30}.\]
\[d_2=d_1-0.20=\boxed{0.10}.\]
\end{frame}

\begin{frame}[t]{Solution 9}
\justifying
\[C=100N(0.30)-100e^{-0.04}N(0.10).\]
\[C=100(0.6179)-100(0.96079)(0.5398)\approx\boxed{9.93}.\]
\end{frame}

\begin{frame}[t]{Solution 10}
\justifying
A sudden jump violates continuous sample paths and challenges constant volatility and normal log-return assumptions. Delta computed immediately before the announcement cannot eliminate jump risk. The hedge may require jump or stochastic-volatility models, wider stress scenarios, and explicit liquidity reserves.
\end{frame}

\begin{frame}[t]{Chapter 9 Takeaways}
\justifying
\begin{itemize}
\item Continuous growth explains the exponential and logarithmic terms.
\item Brownian scaling explains $dW$, quadratic variation, and the It\^o correction.
\item Delta hedging and no arbitrage produce the Black--Scholes PDE.
\item The normal CDF enters when the lognormal payoff integral is evaluated.
\item The model provides a benchmark, while actual risk management must address jumps, changing volatility, costs, and liquidity.
\end{itemize}
\end{frame}

\end{document}
